\subsection{合成律}

定义 1. 设 E 是一个集合。从 E × E 到 E 的一个映射 f 称为 E 上的一个合成律。f 在有序对 (x, y) \ensuremath{\in} E × E 处的值 f(x, y) 称为 x 与 y 在该律下的合成。带有一个合成律的集合称为广群。

x 与 y 的合成通常按确定的顺序写出 x 和 y，并用该律的一个特征符号将它们隔开来表示（这个符号可以约定省略）。最常用的符号是 + 和 ·，通常的约定是后者在需要时可以省略；用这些符号时，x 与 y 的合成分别写成 \(x + y\) 和 x·y 或 xy。用符号 + 表示的律通常称为加法（其合成 \(x + y\) 称为 x 与 y 的和），并说它是加法记法；用符号 · 表示的律通常称为乘法（其合成 x·y = xy 称为 x 与 y 的积），并说它是乘法记法。在本章第1至3节的一般论述中，我们通常用符号 \(\top\) 和 \(\bot\) 来表示任意的合成律。

作为一种语言的滥用，从 \(E \times E\) 的某个子集到 E 的映射有时也称为 E 上并非处处有定义的合成律。

例. (1) 映射 \((\mathbf{X},\mathbf{Y})\mapsto \mathbf{X}\cup \mathbf{Y}\) 和 \((\mathbf{X},\mathbf{Y})\mapsto \mathbf{X}\cap \mathbf{Y}\) 是集合 E 的子集集合上的合成律。

\begin{enumerate}
\def\labelenumi{(\arabic{enumi})}
\setcounter{enumi}{1}
\item
  在自然数集 \(\mathbf{N}\) 上，加法、乘法和乘方都是合成律（在这些律下 \(x \in \mathbf{N}\) 和 \(y \in \mathbf{N}\) 的合成分别记为 \(x + y\) , \(xy\) 或 \(x.y\) 和 \(x^y\) ）（《集合论》第三章第3节第4号）。
\item
  设 E 是一个集合；映射 \((\mathbf{X}, \mathbf{Y}) \mapsto \mathbf{X} \circ \mathbf{Y}\) 是 \(E \times E\) 的子集集合上的一个合成律（《集合论》第二章第3节第3号定义6）；映射 \((f, g) \mapsto f \circ g\) 是从 E 到 E 的映射集合上的一个合成律（《集合论》第二章第5节第2号）。
\item
  设 E 是一个格（集合论，III，§ 1，第 11 号）；若 \(\sup(x, y)\) 表示集合 \(\{x, y\}\) 的最小上界，则映射 \((x, y) \mapsto \sup(x, y)\) 是 E 上的一个合成律。对于最大下界 \(\inf(x, y)\) 也有类似结论。上面的例1是此情形的一个特例，其中 \(\mathfrak{P}(E)\) 按包含关系排序。
\item
  设 \((\mathbf{E}_i)_{i\in I}\) 为一族广群。令 \(\mathsf{T}_i\) 表示 \(\mathbf{E}_i\) 上的合成律。映射
\end{enumerate}

\[
((x _ {i}), (y _ {i})) \mapsto ((x _ {i} \top_ {i} y _ {i}))
\]

是乘积 \(E = \prod_{i \in I} E_i\) 上的一个合成律，称为诸律 \(T_i\) 的乘积。带有该律的集合 E 称为诸广群 \(E_i\) 的乘积广群。特别地，如果所有广群 \(E_i\) 都等于同一个广群 M，则得到由 I 到 M 的映射构成的广群。

设 \((x,y)\mapsto x\top y\) 为集合 E 上的一个合成律。给定 E 的任意两个子集 X、Y，X \(\top\) Y（若此记号不致引起混淆†）表示 E 中诸元素 \(x\top y\) 所成的集合，其中 \(x\in X\) , \(y\in Y\) （换言之，即 X \(\times\) Y 在映射 \((x,y)\mapsto x\top y\) 下的像）。

若 \(a \in \mathbf{E}\) ，我们通常记 \(a \top \mathbf{Y}\) 以代替 \(\{a\} \top \mathbf{Y}\) ，并记 \(\mathbf{X} \top a\) 以代替 \(\mathbf{X} \top \{a\}\) 。映射 \((\mathbf{X}, \mathbf{Y}) \mapsto \mathbf{X} \top \mathbf{Y}\) 是 \(\mathbf{E}\) 的子集所成集合上的一个合成律。

定义2. 设 E 为一个广群，\(\top\) 表示其合成律。E 上的合成律 \((x, y) \mapsto y \top x\) 称为上述合成律的反合成律。带有该律的集合 E 称为 E 的反广群。

设E和E'是两个广群；我们用同一个符号 T 表示它们的律。按照一般定义（《集合论》第四章第1节第5段），一个从E到E'上的双射映射f，使得

\[
f (x \top y) = f (x) \top f (y)\tag{1}
\]

对每个有序对 \((x, y) \in \mathbf{E} \times \mathbf{E}\) 成立，这样的映射称为 \(\mathbf{E}\) 到 \(\mathbf{E}'\) 上的一个同构。若存在 \(\mathbf{E}\) 到 \(\mathbf{E}'\) 上的一个同构，则称 \(\mathbf{E}\) 与 \(\mathbf{E}'\) 同构。

更一般地：

定义3. 一个映射 \(f\) （把 \(\mathbf{E}\) 映入 \(\mathbf{E}'\) ），若关系式(1)对每个有序对 \((x, y) \in \mathbf{E} \times \mathbf{E}\) 成立，则称为 \(\mathbf{E}\) 到 \(\mathbf{E}'\) 的同态或态射；若 \(\mathbf{E} = \mathbf{E}'\) , \(f\) 称为 \(\mathbf{E}\) 的自同态.

广群 E 的恒等映射是一个同态，两个同态的复合也是一个同态。

映射 \(f\) （从 \(E\) 到 \(E'\) ）成为同构的充要条件是：它是一个双射同态，且此时 \(\overline{f}^{-1}\) 是 \(E'\) 到 \(E\) 的一个同构.

